\section{Interval exchange transformations}
\label{sc:IETs}

\subsection{Notations and basic properties} An \emph{interval exchange transformation (IET)} $T$ on a bounded interval $I$ is a piecewise linear bijection of $I$, with finite number of intervals of continuity, on which $T$ acts via translation. {For convenience and without loss of generality, we~assume the interval $I$ to be of the form $[a, b)$, for some $a, b \in \R$, and the IETs to be right-continuous.}

More precisely, there exists $\A$ is an alphabet of $d\in\N$ elements, a~permutation $\pi=\binom{\pi_0}{\pi_1}$ with $\pi_0,\pi_1:\mathcal A\to\{1,\ldots,d\}$ and a collection of left closed and right open subintervals $\{I_{\alpha}\}_{\alpha\in\mathcal A}$ such that and $\bigsqcup_{\alpha\in\A} I_{\alpha}=I$, $T|_{I_{\alpha}}$ acts via translation, and $\pi_0$ and $\pi_1$ describe the order of intervals respectively before and after
action of $T$. It~is easy to check that $T$ preserves Lebesgue measure on $I$ and that the parameters $(\pi,\lambda)$ fully describe the dynamics of $T$, where $\lambda=[\lambda_\alpha]_{\alpha\in\A}:=[|I_{\alpha}|]_{\alpha\in\A}\in \Lambda^{\mathcal A,I}$ is the vector of lengths of intervals $I_{\alpha}$, with $\Lambda^{\mathcal A,I}:=\{\lambda\in\R^{\A}_{+}\mid \sum_{\alpha\in\A}\lambda_{\alpha}=|I|\}$. {Moreover, we~always assume that $\pi$ is \emph{non-reducible}, that is
\[
\pi_1\circ\pi^{-1}_0\{1,\ldots,k\}=\{1,\ldots,k\}\implies k=d.
\]
Otherwise, we~can decompose $T$ into two non-trivial IETs and consider their properties separately. We denote by $S_0^{\A}$ the set of all non-reducible permutations of alphabet~$\A$.
}

For every $\alpha\in\A$, we~denote by $\partial I_{\alpha}$ the left endpoint of $I_{\alpha}$ and by $c_{\alpha}$ its center point. We say that an IET $T$ has a \emph{connection} if there exist $\alpha,\beta\in\A$ with $\pi_0(\beta)\neq 1$, $\pi_1(\alpha)\neq 1$ and $n\in\N_+$ such that
\[
T^{-n}(\partial I_\beta)=\partial I_{\alpha}.
\]
By a connection we often mean the orbit segment $\{T^{-k}(\partial I_\beta)\}_{k=-n,\ldots,0}$.
If such connection exists, we~denote
\[
M(\beta)=M(T,\beta):=\min_{\alpha\in\A}\min\{n\in\N_+\mid T^{-n}(\partial I_\beta)=\partial I_{\alpha}\}.
\]
Otherwise, we~write $M(\beta)=\infty$. Similarly, we~denote
\[
N(\alpha)=N(T,\alpha):=\min_{\beta\in\A}\min\{n\in\N_+\mid T^{n}(\partial I_\alpha)=\partial I_{\beta}\}
\]
and write $N(\alpha)=\infty$ if such connection does not exist. {Note that we always have $T(\partial I_{\pi_1^{-1}(1)})=\partial I_{\pi_0^{-1}(1)}$, a~\emph{trivial} connection. Hence, we~define $M(\pi_0^{-1}(1)):=1$ and $N(\pi_1^{-1}(1)):=1$. {We say that a connection is \emph{irreducible} if it is not a concatenation of two connections, \ie it contains at most two points of the form $\partial I_{\alpha}$, $\alpha\in\A$}.

}

Note that the existence of a non-trivial connection implies that some
non-trivial integer combination of lengths of exchanged intervals is equal to
$0$. Thus, if the length vector is \emph{rationally independent}, i.e., if
$\sum_{\alpha\in\A}r_{\alpha}\lambda_\alpha=0$ for some $(r_{\alpha})_{\alpha
\in \A} \in\mathbb Q^{\A}$ implies that $r_{\alpha}=0$, for all $\alpha\in\A$,
then there cannot be any non-trivial connection. Hence, almost every IET has no non-trivial
connections.

However, in this article we consider the class of \emph{all} ergodic IETs, and, let us recall, the existence of non-trivial connections does not exclude ergodicity. Nevertheless, it~is well-known that if all $\partial I_{\beta}$ are endpoints of connections, then such an IET cannot be ergodic. For the sake of completeness, let us provide a proof of this fact.

\begin{lemma}\label{lem: notallconnections}
Let $T:I\to I$ be an interval exchange transformation given by a permutation $\pi$ and length vector $\lambda$. If for every $\beta\in \A\setminus \{\pi_0^{-1}(1)\}$ we have $M(\beta)<\infty$, then $T$ has only periodic orbits. {More precisely, the base interval $I$ can be decomposed in a finite number of periodic components, given by semi-closed intervals, such that the period is uniform on each of these components.}
\end{lemma}

\begin{proof}
Note that by assumption the set of points
\[
\{T^{n}(\partial I_{\alpha})\mid n\in\Z,\quad\alpha\in \A\}=\{T^{-n}(\partial I_{\alpha})\mid \alpha\in\A\text{ and } 0\le n\le M(\alpha)\}
\]
is finite. Consider the partition given by those points and let $[a,b)$ be an element of this partition.

Note that $T^n$ acts continuously on $[a,b)$ for all $n\in\N$. Indeed, the only possible points of discontinuity of $T$ are $\{\partial I_{\alpha}\}_{\alpha\in\A}$, hence, if for some $n\in\N$ the map $T^n$ did not act continuously on $[a,b)$, there would exist $\beta\in\mathcal A$ such that $\partial I_{\beta}\in T^{n-1}\big([a,b)\big)$, which contradicts the choice of $[a,b)$. In particular, it~follows that $T^n\big([a,b)\big)$ is an interval, for any $n \in \N$.

By Poincaré's recurrence theorem, there exists $N\in\N$ such that
\[
T^{N}\big([a,b)\big)\cap [a,b)\neq \emptyset.
\]
This implies that $T^{N}\big([a,b)\big)=[a,b)$. Indeed, otherwise either $T^{-N}(a)$ or $T^{N}(a)$ belongs to $[a,b)$. Since $a$ is in the orbits of one of the points $\{\partial I_{\alpha}\}_{\alpha\in\A}$, this yields a contradiction.

To sum up, $T^{N}\big([a,b)\big)=[a,b)$ and $T^N$ acts continuously on $[a,b)$. Since $T$ is a piecewise translation, then so is $T^N$. Thus $T^N|_{[a,b)}$ is the identity map on $[a,b)$, which finishes the proof.
\end{proof}

\begin{remark}
By proceeding symmetrically, one can replace in Lemma \ref{lem: notallconnections} the endpoints of connections with their initial points.
\end{remark}

One of the main consequences of the above lemma is the following.

\skpt
\begin{corollary}\label{cor: existenceofinfiniteorbit}
Assume that $T$ is an ergodic IET. Then there exists {$\beta \in\mathcal A \setminus \{\pi_0^{-1}(1)\}$ such that $M(\beta) = +\infty$. }
\end{corollary}

\begin{proof}
Assume, for the sake of contradiction, that $T$ is ergodic but that the conclusion does not hold. Then, by~Lemma \ref{lem: notallconnections}, there exists a non-empty semi-closed interval $[a, b) \subseteq I$ and $N \geq 1$ such that $T^N\mid_{[a, b)}$ is the identity map in $[a, b)$. Therefore, the set $\bigcup_{i = 0}^{N - 1} T^i\big( \big[a, \psfrac{a + b}{2}\big]\big)$ is a non-trivial $T$-invariant set, which contradicts the ergodicity of $T$.
\end{proof}

\subsection{Induced IETs}\label{sc:induced_IETs}
Throughout the proofs of the main results of this paper, we~will often use the first return map of $T$ to a subinterval $J \subseteq I$, {which we denote by $T_J: J \to J$. More precisely, we~define $T_J$ as $x \mto T^{r_J(x)}(x)$, where $r_J: J \to \N$ is given by
\[
 r_J(x):= \min\{n \geq 1 \mid T^n(x) \in J\}.
\]
We sometimes refer to $T_J$ as the \emph{induced map of $T$ to $J$}.

A priori, the map $T_J$ is not necessarily well-defined for \emph{all} points in $J$, although Poincaré's recurrence theorem guarantees that $T_J$ is well-defined in a full Lebesgue measure subset of $J$. However, it~is well known (see, e.g., \cite[\S 3]{veech_interval_1978}) that for any subinterval $J=[a_J,b_J)\subseteq I$ the induced map $T_J$ is an IET of at most $d+2$ intervals, where the possible discontinuities are given by preimages of the discontinuities of $T$ (at most $d - 1$ points) and of the endpoints of $J$ (at most $2$ points, not necessarily disjoint with the previous set).

More precisely, the possible discontinuity points of $T_J$ are given by
\[
\{T^{-m_{J,\alpha}}(\partial I_{\alpha})\}_{\alpha\in \A \setminus \{\pi_0^{-1}(1)\}}, \quad\ m_{J,\alpha}:=\inf\{n\geq 0 \mid T^{-n}(\partial I_{\alpha})\in \mathring{J}\} \text{ \ for } \alpha \in \A,
\]
together with
\[
T^{-m_l}(a_J), \quad m_l :=\inf\{n\ge 0 \mid T^{-n}(a_J)\in \mathring{J}\},
\]
if $a_J$ is different from the left endpoint of $I$, and
\[
T^{-m_r}(b_J), \quad m_r :=\inf\{n\ge 0 \mid T^{-n}(b_J)\in \mathring{J}\},
\]
if $b_J$ is different from the right endpoint of $I$,} where the preimages for
which $m_{J, \alpha}$ (\resp $m_l$ or $m_r$) is $+\infty$ are disregarded.
However, note that if $T$ is minimal, which is the case if $T$ is ergodic (see
Lemma \ref{lem: minimality}), all the above notions are finite.

Moreover, if $T$ is ergodic and $J=[a_J,b_J)\subseteq I$ is chosen so that
\begin{equation}\label{eq: paramofJ}
\begin{gathered}
\text{$a_J=T^{m_0}(\partial I_{\alpha_J})$ and $b_J=T^{n_0}(\partial I_{\beta_J})$, \quad for some $\alpha_J,\beta_J\in\mathcal A$ and $m_0,n_0\in\Z$,} \\
T^{m}(\partial I_{\alpha_J})\notin J, \qquad \text{for any } m \in \{0, \ldots, m_0\} \text{ with } m \neq m_0 \\
T^n(\partial I_{\beta_J})\notin J \qquad \text{for any } n \in \{0, \ldots, n_0\} \text{ with } n \neq n_0,
\end{gathered}
\end{equation}
then the induced map $T_J$ can be seen as an IET of at most $d$ intervals. Indeed, in this case, the discontinuities of $T_J$ belong to the set $
\{T^{-m_{J,\alpha}}(\partial I_{\alpha})\}_{\alpha\in \A\setminus \{\pi_0^{-1}(1)\}}$. Analogously, if $J$ is of the form \eqref{eq: paramofJ} the discontinuities of $T_J^{-1}$ are contained in
\[
\{T^{n_{J,\alpha}}(\partial I_{\alpha})\}_{\alpha\in \A\setminus \{\pi_1^{-1}(1)\}},\!\quad \text{where}\ n_{J,\alpha}:=\min\{n\geq 1\mid T^{n}(\partial I_{\alpha})\in J\} \text{ for any } \alpha \in \A.
\]
If, in addition, $T$ has no non-trivial connections, then the previous two sets have exactly $d - 1$ elements, and $T_J$ can be naturally seen as an IET on $d=\#\A$ intervals and identified with an element $(\pi_J, \lambda_J) \in S_0^\A \times \Lambda^{\A, J}$.

The following simple auxiliary fact tells us what happens if $T$ has non-trivial
connections.

\begin{lemma}\label{lem: alphabetafetrconnections}
Let $T$ be an ergodic interval exchange transformation of $d=\#\A$ intervals and let $J$ be a subinterval of the form \eqref{eq: paramofJ}. Then $T_J$ can be considered as an interval exchange of $d_J$ intervals, where
\[
d_J:= d - \#\{\alpha\in\A\setminus \{\pi_0^{-1}(1)\} \mid m_{J,\alpha}\ge M(\alpha)\}.
\]
In particular, if $J$ does not contain any point from any connection, then $d-d_J$ is equal to the number of non-trivial irreducible connections of $T$.
\end{lemma}

\skpt
\begin{proof}
Since we know that the discontinuities of $T_J$ belong to $\{T^{-m_{J,\alpha}}(\partial I_{\alpha})\}_{\alpha\in \A\setminus \{\pi_0^{-1}(1)\}}$ which has at most $d-1$ elements, to prove that $T_J$ can be seen as an interval exchange of $d_J$ intervals it is sufficient to show that this set has exactly $d_J-1$ elements.

Assume that $\alpha\in\A\setminus \{\pi_0^{-1}(1)\}$ is such that $m_{J,\alpha}\ge M(\alpha)$ and let $\beta\in \A\setminus \{\pi_0^{-1}(1)\}$ be such that $T^{-M(\alpha)}(\partial I_{\alpha})=\partial I_{\beta}$. Then, by~the assumption on $\alpha$, we~have that
\[
T^{-m_{J,\alpha}}(\partial I_{\alpha})=T^{-m_{J,\beta}}(\partial I_{\beta}).
\]
This shows that the connection that ends in the point $\partial I_{\alpha}$ decreases the number of discontinuities of $T_J$ by 1.
To conclude the proof of the first statement it remains to repeat the above reasoning for all $\alpha\in\A$ satisfying $m_{J,\alpha}\ge M(\alpha)$.

To prove the second assertion, first remark that
\[
\#\{\alpha\in\A\setminus \{\pi_0^{-1}(1)\} \mid M(\alpha) < +\infty\}
\]
is equal to the number of non-trivial irreducible connections. Then it suffices to notice that
\[
\#\{\alpha\in\A\setminus \{\pi_0^{-1}(1)\} \mid m_{J,\alpha}\ge M(\alpha)\} = \#\{\alpha\in\A\setminus \{\pi_0^{-1}(1)\} \mid M(\alpha) < +\infty\},
\]
if $J$ does not contain any point from any connection.
\end{proof}

In view of the previous lemma, throughout this work, if $T: I \to I$ is an ergodic IET and $J \subseteq I$ is of the form \eqref{eq: paramofJ}, we~will consider the induced IET $T_J$ as an IET on $d_J$ intervals, where
\begin{align*}
d_J &= 1 + \# \{T^{-m_{J,\alpha}}(\partial I_{\alpha}) \mid \alpha\in \A\setminus \{\pi_0^{-1}(1)\}\\
&= d -\#\{\alpha\in\A\setminus \{\pi_0^{-1}(1)\} \mid m_{J,\alpha}\ge M(\alpha)\} \leq d.
\end{align*}
We will also identify $T_J$ with an element $(\pi_J, \lambda_J) \in S_0^{\A_J} \times \Lambda^{\A_J, J}$ of a possibly smaller alphabet $\A_J$ and denote by $\{I_\gamma^J\}_{\gamma \in \A_J}$ the intervals exchanged by $T_J$.

Let us point out that, in the same way that an IET $T$ with $d$ intervals might have less than $d - 1$ discontinuities, the induced map $T_J$ might have less than $d_J - 1$ discontinuities, that is, some of the points in $
\{T^{-m_{J,\alpha}}(\partial I_{\alpha})\}_{\alpha\in \A\setminus \{\pi_0^{-1}(1)\}}$ might not be real discontinuity points of $T_J$.

Notice that given an ergodic IET $T: I \to I$ and a subinterval $J \subseteq I$ of the form~\eqref{eq: paramofJ}, by~the minimality of $T$ and $T^{-1}$, we~can express $I$ as a disjoint union of the form
\begin{equation}
\label{eq:towers_decomposition}
I = \bigsqcup_{\gamma \in \A_J} \bigsqcup_{i = 0}^{h_\gamma - 1} T^i(I_\gamma^J),
\end{equation}
where, for any $\gamma \in \A_J$, $h_\gamma$ denotes the first return time to $J$ by $T$ of any point in $I_\gamma^J$. We refer to the above decomposition as the \emph{Rokhlin tower decomposition associated with $T_J$ and to $(h_{\gamma})_{\gamma\in \A_J}$ as the corresponding \emph{height} vector}.

We finish this section by recalling a well-known fact, which we prove for completeness.

\begin{lemma}\label{lem: minimality}
If $T: I \to I$ is ergodic with respect to the Lebesgue measure then it is minimal.
\end{lemma}

\begin{proof}
We will show that for every $x,y\in I$ and every $\epsilon>0$ there exists $m\in\N$ such that $|T^{-m}y-x|<2\epsilon$. Take an interval $J:=[x,x+\epsilon)$ and consider the first return map $T_J$. Let $I^J_{\beta}$ be intervals exchanged by $T_J$ and $h_{\beta}$ the corresponding first return times. Since $T$ is ergodic, the set $\tilde I:=\bigcup_{\beta\in \mathcal B}\bigcup_{k=0}^{h_{\beta}-1} T^k(I^J_{\beta})$ is of full Lebesgue measure.

Define $h:=\max\{h_{\beta}\mid\beta\in\mathcal B\}$. Consider the set
\[
C:=\{T^{j}(\partial I_{\alpha})\mid \alpha\in\A,\ j=0,\ldots,h\}.
\]
Pick $0<\delta<\epsilon$ such that $(y,y+\delta]\cap C=\emptyset$. Since $\tilde I$ is of full measure, there exists $\tilde y\in[y,y+\delta]\cap \tilde I$. By the choice of $\delta$ and by the fact that $T$ is right-continuous, the sets $\{T^{-j}[y,\tilde y]\mid j=0,\ldots,h\}$ are a family of pairwise disjoint intervals and $T^{-1}$ acts on each of them by translation. Since $\tilde y\in\tilde I$, there exists $m\le h$ such that $T^{-m}\tilde y\in J$ and $T^{-m}[y,\tilde y]=[T^{m}y,T^{-m}\tilde y]$. Thus
\[
|T^{-m}y,x|<\epsilon+\delta<2\epsilon,
\]
which finishes the proof.
\end{proof}
